% =========================================================================
%  I. Introduction
% =========================================================================
\section{Introduction}

Most classic analog synthesizers create sound with subtractive synthesis.
An oscillator produces a waveform that is rich in overtones, and a filter
removes part of them to shape the timbre \cite{valimaki2007antialiasing}.
The most common starting waveform is the sawtooth, which rises in a
straight line and then drops back at once, and software synthesizers
have to compute it digitally.

A digital signal is a sequence of samples taken at a sampling rate $\fs$,
here 44.1 kHz, and it can only represent frequencies up to half of it,
the Nyquist frequency. A tone with the fundamental frequency $\fo$
consists of harmonics at $\fo$, $2\fo$, $3\fo$ and so on. Because the
sawtooth jumps once per period, its harmonics never end, and their
amplitudes fall only as $1/k$. When the waveform is sampled directly,
every harmonic above the Nyquist frequency folds back into the audible
range. A harmonic at 30 kHz, for example, appears at
14.1 kHz. These aliases do not belong to the harmonic series of the note,
so they sound like noise or out-of-tune tones. An analog recording removes them with a filter before
sampling, but a digital oscillator creates the samples directly, so
aliasing has to be prevented during generation.

Section~\ref{sec:overview} reviews three IEEE journal articles on this
problem, Section~\ref{sec:contribution} compares two cheap methods with a
naive and an ideal oscillator, and Section~\ref{sec:conclusions}
concludes.
